\documentclass[reprint,twocolumn,prb,amsmath,amssymb,superscriptaddress,floatfix]{revtex4-2}

\usepackage{graphicx}
\usepackage{bm}
\usepackage[T1]{fontenc}
% Times-compatible text and math, matching the STIX serif used in the figures,
% so that a symbol in a caption and the same symbol inside a panel agree.
\usepackage{newtxtext}
\usepackage{newtxmath}
\usepackage{booktabs}
\usepackage{xcolor}
\usepackage{xspace}

% Same pastel link palette as the main text.
\definecolor{pastellink}{rgb}{0.35,0.45,0.75}
\definecolor{pastelcite}{rgb}{0.20,0.52,0.45}
\definecolor{pastelurl}{rgb}{0.52,0.36,0.60}

\usepackage[colorlinks=true,
            linkcolor=pastellink,
            citecolor=pastelcite,
            urlcolor=pastelurl,
            filecolor=pastelurl,
            menucolor=pastellink,
            anchorcolor=pastellink,
            linktoc=all,
            breaklinks=true]{hyperref}

% Prefix every PDF destination so that concatenating main.pdf and
% supplement.pdf cannot produce colliding anchor names.
\renewcommand*{\HyperDestNameFilter}[1]{supplement.#1}

\graphicspath{{figures/}}
% ---------------------------------------------------------------------------
% dft_numbers.tex
%
% Every number in the manuscript that comes from the DFT+U campaign is defined
% here and nowhere else, in both the article and the Supplemental Material.
%
%   campaign (a)  results/dftu_row_ordered      row-ordered Li_{1/2}CoO2
%   campaign (b)  results/dftu_order_contrast   row / zigzag / clumped, x = 1/2
%   campaign (c)  results/dftu_x13              sqrt3 x sqrt3 Li_{1/3}CoO2
%   distillation  results/dftu_summary.json
%
% Usage convention: a macro is followed by "~" and then its unit, e.g.
% \comomentRow~$\mu_{\mathrm{B}}$.  "~" is in the xspace exception list.
%
% PROVENANCE NOTE (verification finding B7).  The x = 1/2 site moments, charge
% split and the moment ceiling derived from them are quoted from the ground
% state, afmlo_intra_U4.0 in cell B1, and not from the ferromagnetic solution of
% the seven-atom cell.  \comomentRowRange records the spread across the FM and
% both AFM configurations, which is the honest measure of how well determined
% the moment is.
%
% RETIRED MACROS, and why.  \deltaEOrder and the original reading of
% \dosContrast assumed the row motif lies below the zigzag motif and the ordered
% motif carries the smaller density of states at E_F; neither survived.
% \dosContrast now names the measured U = 0 ratio, which is the falsification
% itself.  \momentCellRow, \enmfmRow, \dosRow, \dosZigzag, \dosClumped,
% \dosXthird and \deltaMdft are retired, the last because the bulk ground states
% are antiferromagnetic and supply no net Delta M.  \jintraRow and \jinterRow
% are retired in favour of \eafmIntra and \eafmInter, which are signed energy
% differences and carry no implicit spin-Hamiltonian convention; \jeffXthird is
% retired in favour of \splitXthird, the raw splitting per cell.
% ---------------------------------------------------------------------------

\newcommand{\Uused}{\ensuremath{4.0}\xspace}

% --- (a) row-ordered Li_{1/2}CoO2, ground state AFM within a Co4+ row ------
\newcommand{\comomentRow}{\ensuremath{0.97}\xspace}       % mu_B on Co4+, ground state
\newcommand{\comomentThreeRow}{\ensuremath{0.00}\xspace}  % mu_B on Co3+, ground state
\newcommand{\comomentRowRange}{\ensuremath{0.97}\ to \ensuremath{1.02}\xspace} % across FM and both AFM
\newcommand{\msatRow}{\ensuremath{0.48}\xspace}           % mu_B per Co if every moment aligned
\newcommand{\dSbound}{\ensuremath{0.23}\xspace}           % k_B per Co, bulk reading of Eq. (S34)
\newcommand{\magorderRow}{antiferromagnetic within a Co\ensuremath{^{4+}} row\xspace}
\newcommand{\efmafmRow}{\ensuremath{28.2}\xspace}         % meV/Co, FM above AFM
\newcommand{\enmRow}{\ensuremath{63.7}\xspace}            % meV/Co, NM above AFM
\newcommand{\gapRow}{\ensuremath{0.70}\xspace}            % eV, AFM ground state (0.698)
\newcommand{\gapRowFM}{\ensuremath{0.43}\xspace}          % eV, FM in the same cell (0.428)
\newcommand{\eafmIntra}{\ensuremath{-56.5}\xspace}        % meV/Co4+, E(AFM) - E(FM), intra-row
\newcommand{\eafmInter}{\ensuremath{+41.7}\xspace}        % meV/Co4+, E(AFM) - E(FM), inter-row
\newcommand{\jintraMag}{\ensuremath{56}\xspace}           % meV, magnitude of the dominant coupling
\newcommand{\dnRow}{\ensuremath{0.20}\xspace}             % e, d-occupation split, ground state
\newcommand{\bondRow}{\ensuremath{0.001}\xspace}          % Angstrom, contrast of site-mean Co-O
\newcommand{\bondSpanRow}{\ensuremath{0.021}\xspace}      % Angstrom, span of individual Co-O bonds

% --- (b) three periodic Li motifs at x = 1/2, 4x2 cell, 8 f.u. ------------
%     Fixed ideal geometry, ferromagnetic alignment imposed.
\newcommand{\deltaEZigzag}{\ensuremath{-8.3}\xspace}       % meV/f.u. vs row, U=4
\newcommand{\deltaEZigzagUzero}{\ensuremath{+8.2}\xspace}  % meV/f.u. vs row, U=0
\newcommand{\deltaEClumped}{\ensuremath{58.1}\xspace}      % meV/f.u. vs row, U=4
\newcommand{\deltaEClumpedUzero}{\ensuremath{92.4}\xspace} % meV/f.u. vs row, U=0
\newcommand{\gapRowMotif}{\ensuremath{0.679}\xspace}       % eV
\newcommand{\gapZigzag}{\ensuremath{0.800}\xspace}         % eV
\newcommand{\gapClumped}{\ensuremath{0.675}\xspace}        % eV
\newcommand{\gapMotifSpread}{\ensuremath{0.13}\xspace}     % eV, 0.800 - 0.675
\newcommand{\dosRowUzero}{\ensuremath{22.5}\xspace}        % states/eV/cell, U=0
\newcommand{\dosContrast}{\ensuremath{1.004}\xspace}       % clumped/row at U=0
\newcommand{\momentRowUzeroA}{\ensuremath{0.3990}\xspace}  % mu_B, row cell at U=0, site 1
\newcommand{\momentRowUzeroB}{\ensuremath{0.3986}\xspace}  % mu_B, row cell at U=0, site 2

% --- (c) sqrt3 x sqrt3 Li_{1/3}CoO2, ground state ferrimagnetic -----------
\newcommand{\comomentXthird}{\ensuremath{1.04}\xspace}     % mu_B on each Co4+
\newcommand{\magorderXthird}{ferrimagnetic, the two Co\ensuremath{^{4+}} moments antiparallel\xspace}
\newcommand{\efmafmXthird}{\ensuremath{12.4}\xspace}       % meV/Co, FM above ferri
\newcommand{\enmXthird}{\ensuremath{230}\xspace}           % meV/Co, NM above ferri
\newcommand{\splitXthird}{\ensuremath{37}\xspace}          % meV per cell, FM above ferri
\newcommand{\gapXthird}{\ensuremath{1.21}\xspace}          % eV, ferri ground state (1.206)
\newcommand{\gapXthirdUfour}{\ensuremath{1.179}\xspace}    % eV, FM at U=4
\newcommand{\gapXthirdUfive}{\ensuremath{1.166}\xspace}    % eV, FM at U=5.5
\newcommand{\dosXthirdNM}{\ensuremath{7.7}\xspace}         % states/eV/cell, NM
\newcommand{\dnXthird}{\ensuremath{0.26}\xspace}           % e, d-occupation split
\newcommand{\bondXthird}{\ensuremath{0.003}\xspace}        % Angstrom, contrast of site-mean Co-O

% --- methodological numbers quoted in the Supplemental Material ------------
\newcommand{\comomentHighU}{\ensuremath{2.28}\xspace}      % mu_B, Co4+ at U=5.5, high spin
\newcommand{\momentFMa}{\ensuremath{1.028}\xspace}         % mu_B, inequivalent FM Co4+, supercell
\newcommand{\momentFMb}{\ensuremath{0.952}\xspace}         % mu_B, the other one
\newcommand{\forcePlateau}{\ensuremath{78}\xspace}         % meV/Angstrom, residual force plateau
\newcommand{\primSign}{\ensuremath{3.1}\xspace}            % meV/Co, primitive-cell FM above NM


\newcommand{\kB}{k_{\mathrm{B}}}
\newcommand{\muBs}{\mu_{\mathrm{B}}}
\newcommand{\Ea}{E_{\mathrm{a}}}
\newcommand{\Ton}{T_{\mathrm{on}}}
\newcommand{\Tcol}{T_{\mathrm{ord}}}

% Supplement equation, figure and table numbering.
\renewcommand{\theequation}{S\arabic{equation}}
\renewcommand{\thefigure}{S\arabic{figure}}
\renewcommand{\thetable}{S\arabic{table}}
\renewcommand{\thesection}{S\arabic{section}}

\begin{document}

\title{Supplemental Material for\\
``Cold Crystallization of the Lithium Sublattice and the Warming-Only
Resistance Anomaly of Delithiated Li$_x$CoO$_2$''}

\author{Santosh Kumar Radha}
\affiliation{Agentfield AI}

\date{\today}

\maketitle

This document gives the derivations that the main text quotes, in the order the main text uses
them. Equations are numbered S1 onward and are referred to from the main text by those numbers.
Section~\ref{sec:symbols} lists every symbol once, with its dimension, so that no symbol carries
two meanings across the two documents.

\section{Mixture formula and the two-probe circuit}
\label{sec:mixture}

\subsection{Bounds on a two-phase channel}

Let the channel have length $L$ and cross-sectional area $A$, and let a fraction $\phi$ of its
volume have transformed from the parent phase P, of resistivity $\rho_{\mathrm{P}}$, into the
ordered phase O, of resistivity $\rho_{\mathrm{O}}$. Two microgeometries bound every other
one~\cite{Wiener1912,Milton2002}. In the series arrangement the transformed material forms slabs that span
the cross section, so the resistances add,
\begin{equation}
  R_{\mathrm{series}} = \frac{L}{A}\bigl[(1-\phi)\rho_{\mathrm{P}}+\phi\rho_{\mathrm{O}}\bigr].
  \label{eq:sm-series}
\end{equation}
With $R_{\mathrm{P}}\equiv\rho_{\mathrm{P}}L/A$ the resistance the same channel would have with
nothing transformed, and with the resistivity ratio
\begin{equation}
  s \equiv \frac{\rho_{\mathrm{O}}}{\rho_{\mathrm{P}}},
  \label{eq:sm-sdef}
\end{equation}
Eq.~\eqref{eq:sm-series} gives the excess exactly,
\begin{equation}
  \frac{\Delta R}{R_{\mathrm{P}}}
  \equiv \frac{R_{\mathrm{series}}-R_{\mathrm{P}}}{R_{\mathrm{P}}}
  = \phi\,(s-1).
  \label{eq:sm-seriesexcess}
\end{equation}
In the parallel arrangement the transformed material forms filaments running the length of the
channel, so the conductances add, and
\begin{equation}
  \frac{\Delta R}{R_{\mathrm{P}}}
  = \frac{\phi\,(1-1/s)}{1-\phi\,(1-1/s)}
  \;\xrightarrow[\ \phi(1-1/s)\ll1\ ]{}\;
  \frac{\phi\,(s-1)}{s}.
  \label{eq:sm-parallel}
\end{equation}
Both bounds have the form quoted as Eq.~(2) of the main text,
$\Delta R/R_{\mathrm{P}}=c\,\phi\,(s-1)$, with the geometry factor $c=1$ for the series limit
and $c=1/s$ for the parallel limit, so that $c\in(0,1]$ whenever the transformed phase is the
more resistive one. Both expressions vanish at $\phi=0$ and at $s=1$, as they must, and both are
monotonic in $\phi$ and in $s$. The quantity the experiment supplies is the product
$c\,\phi\,(s-1)$, which the zero-field maximum places at approximately unity; the individual
factors are not separately determined by a single trace.

\subsection{What the two-probe geometry allows}

The measured resistance is the sum in Eq.~(1) of the main text,
$R_{\mathrm{meas}}=R_{\mathrm{c}}+R_{\mathrm{ch}}$, where $R_{\mathrm{c}}$ collects the two
Schottky junctions evaluated on the quasilinear branch of their current-voltage characteristic
and $R_{\mathrm{ch}}$ is the channel. Writing the warming and cooling traces separately and
subtracting,
\begin{equation}
  \Delta R = \bigl[R_{\mathrm{c}}^{\mathrm{w}}-R_{\mathrm{c}}^{\mathrm{c}}\bigr]
           + \bigl[R_{\mathrm{ch}}^{\mathrm{w}}-R_{\mathrm{ch}}^{\mathrm{c}}\bigr],
  \label{eq:sm-circuit}
\end{equation}
so that any contact resistance that depends on temperature but not on thermal history cancels
identically. What does not cancel is a contact whose barrier height itself depends on the
lithium arrangement beneath the electrode, since that arrangement carries the same history
dependence as the channel. Equation~\eqref{eq:sm-circuit} is therefore an upper bound on the
channel contribution, and every quantitative statement in the main text is made about
$\Delta R$ rather than about a resistivity. The two experiments that separate the terms are
channel-length scaling, under which the second bracket grows linearly with $L$ while the first
does not, and a four-probe measurement, which removes the first bracket outright.

\section{Unfreezing on the warming branch}
\label{sec:unfreezing}

\subsection{Growth from a frozen nucleus population}

Start from the growth velocity of Eq.~(7) of the main text, with the sample quenched into the
untransformed state and the ordered phase thermodynamically favored below $\Tcol$:
\begin{equation}
  v(T) = v_{0}\,e^{-\Ea/\kB T}\,q(T),
  \qquad
  q(T)\equiv 1-e^{\Delta f/\kB T},
  \label{eq:sm-master}
\end{equation}
in which $v_{0}=\nu a$ is the velocity a domain boundary would have if every attempt succeeded,
$a=0.28$~nm being the lithium-lithium spacing in the plane, and $\Delta f=f_{\mathrm{O}}-
f_{\mathrm{P}}$ is the free-energy difference per cobalt, negative below $\Tcol$. The mobility
factor counts the attempts that clear the migration barrier and the thermodynamic factor $q$
removes the reverse hops, so that $v$ vanishes at $\Tcol$ as it must and approaches
$v_{0}e^{-\Ea/\kB T}$ only at undercoolings large enough that $|\Delta f|$ exceeds $\kB T$.
Close to the boundary $|\Delta f|\simeq\Delta S\,(\Tcol-T)$, so that
$q=1-e^{-\sigma(\Tcol-T)/T}$ with $\sigma=\Delta S/\kB$ per cobalt.

Nucleation is complete before the warming sweep begins, since the cooldown delivers the sample
to base temperature already seeded, so the transformation is site saturated: each of the $n_{0}$
nuclei per unit area grows as a disk of the same radius, and the
Kolmogorov--Johnson--Mehl--Avrami construction of Sec.~\ref{sec:kjma} collapses to
\begin{equation}
  \phi(T) = 1-\exp\bigl[-\pi n_{0}R_{\mathrm{w}}(T)^{2}\bigr],
  \label{eq:sm-formal}
\end{equation}
with $R_{\mathrm{w}}$ the radius accumulated since the start of the sweep. On a linear sweep
$T(t)=T_{0}+rt$ with $r>0$ the substitution $dt'=dT'/r$ turns that radius into a temperature
integral,
\begin{equation}
\begin{aligned}
  R_{\mathrm{w}}(T) &= \frac{1}{r}\int_{T_{0}}^{T}v(T')\,dT' \;\equiv\; a\,\Theta(T),\\[2pt]
  \Theta(T) &= \frac{\nu}{r}\int_{T_{0}}^{T}e^{-\Ea/\kB T'}\,q(T')\,dT'.
\end{aligned}
  \label{eq:sm-theta-def}
\end{equation}
The integrand of $\Theta$ has dimensions of inverse time, the integration variable has
dimensions of temperature, and $r$ has dimensions of temperature per time, so $\Theta$ is
dimensionless, as it must be if $R_{\mathrm{w}}$ is to carry the dimensions of $a$. Its
physical meaning is the number of successful rearrangement attempts made between $T_{0}$ and
$T$, net of the reverse hops that $q$ subtracts, and $R_{\mathrm{w}}=a\Theta$ is the distance a
domain boundary advances over the same interval.

\subsection{Asymptotic evaluation}

Change variable to $y=\Ea/\kB T'$, so that $T'=\Ea/\kB y$ and $dT'=-(\Ea/\kB)y^{-2}dy$.
Equation~\eqref{eq:sm-theta-def} becomes
\begin{equation}
  \Theta(T) = \frac{\nu\Ea}{r\kB}\int_{y}^{y_{0}}
  q\!\left(\frac{\Ea}{\kB y'}\right)\frac{e^{-y'}}{y'^{2}}\,dy',
  \qquad
  y=\frac{\Ea}{\kB T},\ \ y_{0}=\frac{\Ea}{\kB T_{0}}.
  \label{eq:sm-beta}
\end{equation}
Because $y_{0}>y$ and the integrand falls exponentially, the lower limit of the sweep
contributes a fraction $e^{-(y_{0}-y)}(y/y_{0})^{2}q(T_{0})/q(T)$ of the total. For the
parameters of interest, $\Ea=0.387$~eV, $T=146$~K and $T_{0}=80$~K, that fraction is
$7.5\times10^{-12}$, so the lower limit may be replaced by infinity. The thermodynamic factor
varies slowly on the scale over which the exponential does, and is carried to first order by
writing the integrand as $e^{h}$ with $h=\ln q-\Ea/\kB T'$; the pure exponential part has the
standard asymptotic expansion
\begin{equation}
  \int_{y}^{\infty}\frac{e^{-y'}}{y'^{2}}\,dy'
  = \frac{e^{-y}}{y^{2}}\left[1-\frac{2}{y}+O(y^{-2})\right],
  \label{eq:sm-asymptotic}
\end{equation}
obtained by integrating by parts once and bounding the remainder. Substituting
Eq.~\eqref{eq:sm-asymptotic} into Eq.~\eqref{eq:sm-beta}, restoring $y=\Ea/\kB T$ and keeping
the first derivative of $q$ gives Eq.~(4) of the main text together with both of its leading
corrections,
\begin{equation}
  \Theta(T) = \frac{\nu}{r}\,\frac{\kB T^{2}}{\Ea}\,q(T)\,e^{-\Ea/\kB T}
  \left[1-\frac{2\kB T}{\Ea}-\frac{\kB T^{2}}{\Ea}\frac{d\ln q}{dT}+\cdots\right].
  \label{eq:sm-theta}
\end{equation}
The two corrections are of the same size and of opposite sign, and neither may be discarded on
the ground that the other is small. At
$T=146$~K and $\Ea=0.387$~eV the expansion parameter is $\kB T/\Ea=1/30.8$, so the Laplace
correction $2\kB T/\Ea$ is $6.5$ percent; at $\sigma=\dSbound$ and $\Tcol=240$~K the
thermodynamic factor is $q(146\ \mathrm{K})=0.138$ and $d\ln q/dT=-0.0162$~K$^{-1}$, so the
second correction is $7.7$ percent with the opposite sign. The two nearly cancel in the value
of $\Theta$, which is why the leading form suffices for the barrier, but they do not cancel in
its logarithmic derivative, which is what the rate shift below
measures, and they are carried separately from here on. The prefactor
$\kB T^{2}/\Ea$ has dimensions of temperature, so $\nu\kB T^{2}/r\Ea$ is dimensionless.

The value of $q$ matters as much as its derivative. Retaining it rather than setting it to
unity multiplies $\Theta$ by $0.138$ at the onset, which at fixed $n_{0}$ lowers the barrier
the onset returns by $\kB\Ton\ln(1/q)=0.025$~eV and at fixed barrier raises the nucleus density
the onset demands by $q^{-2}$, a factor of $53$. No step of the argument may treat the driving
force as saturated at 146~K, since the undercooling there is only $94$~K against a latent
entropy of a fraction of $\kB$.

Two limits confirm Eq.~\eqref{eq:sm-theta}. As $r\to0$ the accumulated radius diverges and
$\phi\to1$ at every temperature at which $v$ is nonzero, which is the statement that an
infinitely slow sweep is an equilibrium measurement. As $r\to\infty$ the radius vanishes and the
sample never transforms, which is the statement that an infinitely fast sweep quenches.

\subsection{The onset and the barrier}

Define the onset as the rise midpoint, $\phi(\Ton)=1/2$, which by Eq.~\eqref{eq:sm-formal}
requires
\begin{equation}
  \pi n_{0}R_{\mathrm{w}}(\Ton)^{2}=\ln 2,
  \qquad
  R_{\mathrm{w}}(\Ton)=\sqrt{\frac{\ln2}{\pi n_{0}}} .
  \label{eq:sm-onsetdef}
\end{equation}
The onset therefore fixes not a barrier but a pair, since the same radius can be reached with a
larger barrier and a denser seeding or with a smaller barrier and a sparser one; at
$n_{0}=10^{-2}$~nm$^{-2}$, one nucleus per hundred square nanometres, the radius the sweep must
deliver is 4.70~nm. Inserting Eq.~\eqref{eq:sm-theta} at leading order, taking logarithms and
solving for $\Ea$ gives Eq.~(5) of the main text,
\begin{equation}
  \Ea = \kB\Ton\,\ln\!\left[\frac{\nu\,a\,\kB\Ton^{2}\,q(\Ton)}
                                 {\Ea\,r\,R_{\mathrm{w}}(\Ton)}\right].
  \label{eq:sm-Ea}
\end{equation}
The argument of the logarithm is dimensionless: $\nu$ carries s$^{-1}$, $a$ and
$R_{\mathrm{w}}$ carry length, $\kB\Ton^{2}/\Ea$ carries K, $q$ is dimensionless and $r$
carries K\,s$^{-1}$. Equation~\eqref{eq:sm-Ea} is implicit in $\Ea$ but converges under direct
iteration because $\Ea$ appears inside a logarithm. Starting from $\Ea=0.40$~eV with
$\Ton=146$~K, $\nu=10^{13}$~s$^{-1}$, $r=1$~K/min and $n_{0}=10^{-2}$~nm$^{-2}$ the successive
iterates are 0.3869, 0.3873, 0.3873~eV, so three passes suffice, and numerical integration of
Eq.~\eqref{eq:sm-theta-def} without any expansion returns 0.3874~eV, the agreement to one part
in four thousand being what remains of the two corrections in Eq.~\eqref{eq:sm-theta} once they
have cancelled against one another.

The error budget follows from differentiating Eq.~\eqref{eq:sm-Ea}. Writing
$\varepsilon=\Ea$, $a_{0}=\kB\Ton$ and $C=\nu a\kB\Ton^{2}q(\Ton)/(r R_{\mathrm{w}})$, so that
$\varepsilon=a_{0}\ln(C/\varepsilon)$,
\begin{equation}
  \frac{\partial\varepsilon}{\partial\ln C}=\frac{a_{0}\,\varepsilon}{\varepsilon+a_{0}},
  \qquad
  \frac{\partial\varepsilon}{\partial\Ton}
  =\frac{\varepsilon/\Ton+2\kB+\kB\Ton\,d\ln q/dT}{1+a_{0}/\varepsilon}.
  \label{eq:sm-sensitivity}
\end{equation}
With $\varepsilon=0.387$~eV and $a_{0}=\kB\Ton=12.58$~meV, the first derivative is 12.19~meV
per $e$-fold of $C$, that is 28~meV per decade; since $C\propto\nu$ and $C\propto1/r$, a decade
of attempt frequency and a decade of sweep rate each move $\Ea$ by 0.028~eV, in opposite
directions, and since $C\propto R_{\mathrm{w}}^{-1}\propto n_{0}^{1/2}$, a decade of nucleus
density moves it by half as much, 0.014~eV. The second derivative is 2.54~meV/K, whose
reciprocal is the 393~K/eV that the width analysis below uses, so the $\pm7$~K uncertainty on
the onset contributes $\pm0.018$~eV, and carrying the latent entropy over 0.10 to 0.69~$\kB$ per cobalt
through $q(\Ton)$ contributes $-0.010$ to $+0.012$~eV. Added in quadrature these give
0.047~eV, quoted as $\Ea=0.39\pm0.05$~eV. Solving Eq.~\eqref{eq:sm-Ea} at the corners of the
plausible ranges $\nu\in[10^{11},10^{14}]$~s$^{-1}$, $r\in[0.5,5]$~K/min and
$n_{0}\in[10^{-4},10^{-1}]$~nm$^{-2}$ gives 0.28~eV at the low corner and 0.44~eV at the high
corner, widening to the interval 0.27 to 0.46~eV once the onset and latent-entropy ranges are
carried to their corners as well.

The degeneracy between the barrier and the nucleus density is exact enough to be written down.
Holding $\Ton$ at 146~K and everything else at the values above,
$\Ea=0.387+0.014\log_{10}(n_{0}/10^{-2}\ \mathrm{nm}^{-2})$~eV, so the plausible band
$n_{0}=10^{-4}$ to $10^{-1}$~nm$^{-2}$, corresponding to nucleus spacings of 100 down to 3~nm,
maps onto $\Ea=0.36$ to 0.40~eV. Neither number is measured on its own; what the onset measures
is the combination, and the rate series of the next subsection is what separates the barrier
from the seeding.

\subsection{The rate shift and its independence of the Avrami exponent}

Take the logarithm of the onset condition Eq.~\eqref{eq:sm-onsetdef}, writing
$R_{\mathrm{w}}=a\Theta$ and using Eq.~\eqref{eq:sm-theta} at leading order,
\begin{equation}
  \ln(\pi n_{0}a^{2})
  +2\left[\ln\frac{\nu}{r}+\ln\frac{\kB\Ton^{2}q(\Ton)}{\Ea}-\frac{\Ea}{\kB\Ton}\right]
  =\ln\ln2 .
  \label{eq:sm-logonset}
\end{equation}
Differentiate Eq.~\eqref{eq:sm-logonset} with respect to $\ln r$ at fixed $\Ea$, $\nu$ and
$n_{0}$. The factor of two divides out; inside the bracket the rate term gives $-1$, the
prefactor gives $2\,d\Ton/(\Ton\,d\ln r)$, the thermodynamic factor gives
$(d\ln q/dT)\,d\Ton/d\ln r$ and the exponential gives
$+\Ea\,d\Ton/(\kB\Ton^{2}\,d\ln r)$, so
\begin{equation}
  \frac{d\Ton}{d\ln r}
  = \left[\frac{2}{\Ton}+\frac{\Ea}{\kB\Ton^{2}}+\frac{d\ln q}{dT}\right]^{-1},
  \label{eq:sm-rateshift}
\end{equation}
which is Eq.~(6) of the main text and reduces, when the driving force is treated as saturated
and $d\ln q/dT$ is dropped, to the Kissinger relation $\kB\Ton^{2}/(\Ea+2\kB\Ton)$ for a
first-order rate process~\cite{Kissinger1957}. Numerically, at $\Ton=146$~K the three terms in
the bracket are $0.0137$, $0.2107$ and $-0.0162$~K$^{-1}$, giving $4.80$~K per $e$-fold, or
$11.1$~K per decade of warming rate; the thermodynamic factor is what carries the last of the
three, and dropping it would return $10.3$~K per decade. Across the barrier range that the
onset admits at $\Ton=146$~K the same expression gives $11.6$, $11.1$ and $10.7$~K per decade
at $\Ea=0.37$, $0.387$ and $0.40$~eV, and integrating Eq.~\eqref{eq:sm-theta-def} numerically
over a decade of rate from 0.5 to 5~K/min gives $11.4$~K. The prediction quoted in the main
text is therefore $11\pm1$~K per decade, the spread being the width of the admissible barrier
range rather than a fitted uncertainty.

Suppose now that the transformed fraction is $\phi=1-\exp(-A_{0}\Theta^{n})$ for a constant
$A_{0}$ and exponent $n$, which covers continuous nucleation with growth ($n=3$ in two
dimensions) and growth from a frozen nucleus population ($n=2$, with $A_{0}=\pi n_{0}a^{2}$),
as well as the simple relaxation $n=1$. The onset condition becomes $A_{0}\Theta^{n}=\ln2$,
whose logarithm is
\begin{equation}
  n\left[\ln\frac{\nu}{r}+\ln\frac{\kB\Ton^{2}q(\Ton)}{\Ea}-\frac{\Ea}{\kB\Ton}\right]
  = \ln\frac{\ln2}{A_{0}}.
  \label{eq:sm-avrami-onset}
\end{equation}
Differentiating with respect to $\ln r$ divides out the overall factor $n$ and returns
Eq.~\eqref{eq:sm-rateshift} unchanged, the constant $A_{0}$ dropping out entirely; direct
numerical solution of Eq.~\eqref{eq:sm-formal} at $n=1$ and $n=2$ confirms the cancellation to
four figures. The onset shift per decade of warming rate is therefore a property of the barrier
and the onset temperature alone, and it does not test which transformation mechanism operates.
The exponent instead controls the width of the rise, which is what the next subsection
quantifies, and it shifts the extracted $\Ea$ only through the constant $A_{0}$ inside a
logarithm, that is by the 0.014~eV per decade of $n_{0}$ already budgeted.

The cancellation carries one condition, and the rate series has to be run so as to respect it.
Equation~\eqref{eq:sm-avrami-onset} treats $A_{0}$, and with it $n_{0}$, as a constant of the
protocol, which it is only if the cooldown that seeds the sample is held fixed while the
warming rate is varied. If the nuclei are thermally seeded then their density falls as the
cooldown is hurried, $n_{0}\propto1/|r_{\mathrm{c}}|$, and differentiating
Eq.~\eqref{eq:sm-logonset} with respect to $\ln|r_{\mathrm{c}}|$ instead gives half the shift
of Eq.~\eqref{eq:sm-rateshift}, that is $+5.5$~K per decade of cooling rate, with the same sign,
since a sparser seeding forces the sweep to grow larger domains before the midpoint is reached.
A protocol that slaves the cooldown to the warming rate therefore returns about $16.5$~K per
decade, half as much again as the growth shift alone, and the two contributions cannot be
separated after the fact. The measurement that tests Eq.~\eqref{eq:sm-rateshift} is a warming
series taken at a single fixed cooling rate.

\subsection{Width of the rise}

The bracket of Eq.~\eqref{eq:sm-rateshift} is $d\ln\Theta/dT$, so the extended area
$A_{\mathrm{ex}}=\pi n_{0}a^{2}\Theta^{2}$ that sits in the exponent of
Eq.~\eqref{eq:sm-formal} grows by a factor of ten every
\begin{equation}
  \delta T_{10\times}=\frac{\ln10}{2}\left(\frac{d\ln\Theta}{dT}\right)^{-1}=5.5\ \mathrm{K}
  \label{eq:sm-decade}
\end{equation}
at $\Ton=146$~K and $\Ea=0.387$~eV, half the 11.1~K that $\Theta$ itself takes. The
transformed fraction passes from 10 to 90 percent as $A_{\mathrm{ex}}$ runs from
$-\ln0.9=0.105$ to $\ln10=2.303$, a factor of 21.9, so
\begin{equation}
  \Delta T_{10\text{--}90}=\delta T_{10\times}\log_{10}(21.9)=7.4\ \mathrm{K},
  \label{eq:sm-width}
\end{equation}
which becomes 7.3~K when Eq.~\eqref{eq:sm-theta-def} is integrated across the rise rather than
expanded about the midpoint, and that is the number carried forward; the corresponding width
for a general exponent is larger by the factor $2/n$, so continuous nucleation would be
narrower still and simple relaxation twice as broad. The observed width, $19^{+6}_{-9}$~K,
exceeds all of them. A spread of barriers $\delta\Ea$ broadens the onset by
$\delta\Ton=\delta\Ea\,(\partial\Ea/\partial\Ton)^{-1}$, the second factor being the 393~K/eV
that Eq.~\eqref{eq:sm-sensitivity} returns, so
$\delta\Ea=0.045$~eV contributes 17.5~K, which added in quadrature to
Eq.~\eqref{eq:sm-width} gives 19.0~K. A barrier spread of that size is not an adjustable
quantity here: the tetrahedral-site divacancy barrier computed for Li$_{x}$CoO$_{2}$ scatters by
$\pm0.04$~eV across the local lithium environments of a single
composition~\cite{VanderVen2001b}, so the width of the rise and the environment dependence of
the barrier are one statement rather than two. The width is accordingly consistent with growth
from a frozen nucleus population carrying the environment scatter the first-principles barriers
already predict, and is not consistent with a sharply cooperative transformation.

\section{Nucleation, growth, and the branch asymmetry}
\label{sec:kjma}

\subsection{The two-dimensional nucleation barrier}

A circular domain of ordered phase of radius $R$ within the lithium plane costs line tension and
gains bulk free energy~\cite{Turnbull1949,Turnbull1950},
\begin{equation}
  \Delta G(R) = -\pi R^{2}\,|\Delta f_{A}| + 2\pi R\,\lambda ,
  \label{eq:sm-dg}
\end{equation}
with $\Delta f_{A}$ the free-energy difference per unit area, negative below $\Tcol$, and
$\lambda$ the domain-wall line tension, of dimension energy per length. Setting
$d\Delta G/dR=0$ gives the critical radius and the barrier,
\begin{equation}
  R^{*}=\frac{\lambda}{|\Delta f_{A}|},
  \qquad
  W^{*}=\Delta G(R^{*})=\frac{\pi\lambda^{2}}{|\Delta f_{A}|},
  \label{eq:sm-wstar}
\end{equation}
whose dimensions are (energy/length)$^{2}$ divided by (energy/area), that is energy. Close to
the boundary the driving force is linear in the undercooling,
$|\Delta f_{A}|\simeq\Delta S_{A}(\Tcol-T)$ with $\Delta S_{A}$ the latent entropy per unit
area, so
\begin{equation}
  \frac{W^{*}}{\kB T}\simeq\frac{\pi\lambda^{2}}{\kB\,\Delta S_{A}}\,
  \frac{1}{T\,(\Tcol-T)},
  \label{eq:sm-wkt}
\end{equation}
which diverges as $T\to\Tcol^{-}$ and falls monotonically as the undercooling grows.

\subsection{Transformed fraction}

With nucleation rate per unit area $I_{\mathrm{n}}$ and growth velocity $v$ as in Eq.~(7) of the
main text, the extended area fraction of the
Kolmogorov--Johnson--Mehl--Avrami construction~\cite{Kolmogorov1937,Johnson1939,Avrami1939,Avrami1940} is
\begin{equation}
  A_{\mathrm{ex}}(t)=\pi\int_{0}^{t}I_{\mathrm{n}}(t')
  \left[\int_{t'}^{t}v(t'')\,dt''\right]^{2}dt',
  \label{eq:sm-aex}
\end{equation}
in which each nucleus born at $t'$ contributes the area of the disk it would have grown by time
$t$ had it met no other domain, and the transformed fraction is
\begin{equation}
  \phi(t)=1-\exp\bigl[-A_{\mathrm{ex}}(t)\bigr].
  \label{eq:sm-kjma}
\end{equation}
The dimensions check: $I_{\mathrm{n}}$ carries (area\,time)$^{-1}$, the bracket carries length,
and the outer $dt'$ carries time, so $A_{\mathrm{ex}}$ is dimensionless. For constant
$I_{\mathrm{n}}$ and $v$, Eq.~\eqref{eq:sm-aex} reduces to
$A_{\mathrm{ex}}=\pi I_{\mathrm{n}}v^{2}t^{3}/3$, the familiar two-dimensional Avrami exponent
of three. If instead nucleation is complete before growth begins, with a frozen areal density
$n_{0}$ of nuclei, then
\begin{equation}
  A_{\mathrm{ex}}(t)=\pi n_{0}\left[\int_{0}^{t}v\,dt'\right]^{2}
  =\pi n_{0}R_{\mathrm{w}}^{2},
  \label{eq:sm-saturated}
\end{equation}
which is Eq.~\eqref{eq:sm-formal} and the $n=2$ member of the $A_{0}\Theta^{n}$ family analyzed
in Sec.~\ref{sec:unfreezing}, with $A_{0}=\pi n_{0}a^{2}$; the onset shift is then
Eq.~\eqref{eq:sm-rateshift} and only the width changes.

Which of the two limits applies is decided by the nucleation rate, and the point that governs
the whole of this section is that the nucleation rate carries the same mobility factor as the
growth velocity,
\begin{equation}
  I_{\mathrm{n}}(T) = I_{0}\,e^{-[\Ea+W^{*}(T)]/\kB T},
  \label{eq:sm-inuc}
\end{equation}
with $I_{0}=\nu/a^{2}$ for homogeneous nucleation on the lithium sublattice and $I_{0}=X\nu$
for heterogeneous nucleation on a template of areal density $X$. Only $W^{*}$ distinguishes
Eq.~\eqref{eq:sm-inuc} from Eq.~\eqref{eq:sm-master}; the ionic mobility that has to unfreeze
before a domain can grow has equally to unfreeze before a domain can be born.

\subsection{Why cooling and warming differ}

The two exponentials in Eqs.~\eqref{eq:sm-master} and \eqref{eq:sm-inuc} move in opposite
directions. By Eq.~\eqref{eq:sm-wkt} the barrier $W^{*}$ is prohibitive just below $\Tcol$ and
falls monotonically as the undercooling grows, while the mobility factor common to both rates
collapses as the temperature falls. Their product has a single maximum, the nose of the
time-temperature-transformation curve~\cite{Uhlmann1972}, which for every line tension considered
here sits between 178 and 211~K. What it does not have is a second maximum at large
undercooling: the product
$T(\Tcol-T)$ in the denominator of Eq.~\eqref{eq:sm-wkt} is itself largest at $\Tcol/2=120$~K,
so below that temperature $W^{*}/\kB T$ rises again as the sample cools and joins the mobility
factor rather than opposing it, $dI_{\mathrm{n}}/dT>0$ holds throughout, and the bottom of the
sweep is where nucleation is least likely rather than most. At 80~K the mobility factor alone is
$e^{-56}$, and a nucleation event there is as forbidden as a growth event.

The requirement that the smooth cooling trace places on the seeding can therefore be stated as
a birth temperature. Smoothness demands that the extended area accumulated on the way down stay
below some $A_{\mathrm{ex}}^{\max}$, which we take as $10^{-2}$, while the warming branch
demands $\pi n_{0}R_{\mathrm{w}}(\Ton)^{2}=\ln2$; since both are quadratic in the radius, any
nucleus present at base temperature must have grown on the way down to no more than
$(A_{\mathrm{ex}}^{\max}/\ln2)^{1/2}=0.12$ of the radius the warming sweep will give it. The
growth accumulated from a birth temperature $T'$ downward is the same Laplace integral as
Eq.~\eqref{eq:sm-theta}, evaluated at $T'$ instead of $\Ton$, so the ratio of the two radii is
$e^{-(\Ea/\kB)(1/T'-1/\Ton)}$ up to algebraic prefactors of order unity, and the condition is
\begin{equation}
  \frac{\Ea}{\kB}\left(\frac{1}{T'}-\frac{1}{\Ton}\right)
  \;\geq\; \frac{1}{2}\ln\frac{\ln2}{A_{\mathrm{ex}}^{\max}} .
  \label{eq:sm-nose}
\end{equation}
The right-hand side is 2.12, and the left-hand side is steep, so the ceiling it sets is almost
independent of the barrier: $T'\leq136.2$, 136.6, 136.9 and 137.8~K at $\Ea=0.37$, 0.387, 0.40
and 0.45~eV. Every nucleus the sample carries to base temperature must have been born below
$\Ton-9$~K. A nucleus born at the nose, by contrast, grows between 0.6~nm and 40~$\mu$m over
the remainder of a 1~K/min cooldown against the 1.5 to 47~nm that Eq.~\eqref{eq:sm-onsetdef}
asks of it, so a single such nucleus per flake transforms the sample on the way down and
destroys the branch asymmetry.

Classical nucleation cannot deliver seeds under that ceiling, and the arithmetic says so
without appeal. Pushing the maximum of Eq.~\eqref{eq:sm-inuc} below 137~K while $W^{*}$ retains
the $1/(\Tcol-T)$ form of Eq.~\eqref{eq:sm-wkt} requires $|d(W^{*}/\kB T)/dT|$ to exceed
$\Ea/\kB T^{2}=0.21$~K$^{-1}$ there, hence $W^{*}/\kB T\simeq100$, and the rate is suppressed
by $e^{-100}$. Put as a constraint on the one free quantity in Eq.~\eqref{eq:sm-wkt}, the line
tension is overdetermined: keeping the cooling trace smooth requires $\lambda$ to be at least
about 53 to 59~meV/nm over the admissible barriers, while delivering the seeding density that
Eq.~\eqref{eq:sm-onsetdef} needs requires it to be at most 38 to 48~meV/nm. Scanning the whole
of the plausible parameter space, $\Ea$ from 0.34 to 0.45~eV, $\nu$ from $10^{11}$ to
$10^{14}$~s$^{-1}$, $\sigma$ from 0.10 to 0.69, $\lambda$ over four decades and the cooling rate
from 1 to 60~K/min, the two requirements never meet: the shortfall in nucleus density is five
to eight decades at 1~K/min and never falls below 1.9 decades anywhere in the space, and the
corner that comes closest does so only by delivering almost no nuclei at all. Heterogeneous
nucleation does not rescue this. A potency factor $\psi$ that replaces $W^{*}$ by $\psi W^{*}$
is exactly equivalent to replacing $\lambda$ by $\lambda\sqrt{\psi}$, so it moves within the
same one-parameter family already scanned, and it does so with a prefactor $X\nu$ smaller than
$\nu/a^{2}$ by some three orders of magnitude, which costs a further one to one and a half
decades. In the opposite limit, where the template removes the barrier altogether, the sites
activate as soon as the sample crosses $\Tcol$, where Eq.~\eqref{eq:sm-tau} gives a relaxation
time of $3\times10^{-4}$~s, and the flake transforms on the way down.

What the argument leaves is an assumption rather than a derived result.
The seeds are taken to be athermal: a quenched population of defect-templated embryos, patches
of ordered lithium pinned to structural imperfections of the flake, which are present at all
temperatures and go supercritical not by a thermal fluctuation but as soon as the shrinking
critical radius of Eq.~\eqref{eq:sm-wstar} falls to their own size. That crossing carries no
mobility factor and therefore no $e^{-\Ea/\kB T}$ suppression, which is precisely the functional
form the classical estimate cannot supply. At 137~K the critical radius is
$R^{*}=\lambda/|\Delta f_{A}|=0.67$, 1.01 and 1.35~nm for $\lambda=20$, 30 and 40~meV/nm, that
is patches of 21, 47 and 83 cobalt sites, so an embryo population whose typical member is of
that size activates in the right window. The assumption carries a condition of its own: the
embryo-size distribution has to be cut off above about 1.3~nm, because a single
2~nm embryo anywhere in the flake goes supercritical near 200~K and sweeps 10~$\mu$m before the
cooldown ends. A cutoff of that kind is natural for patches templated by a fixed defect
population and unnatural for thermal fluctuations of the parent phase, which is the sense in
which the assumption is a physical one rather than a fitted one.

The assumption is not free, because it makes two predictions that the thermal alternative does
not. If the seeds are athermal then $n_{0}$ is a property of the flake and not of the protocol,
so $\Ton$ is independent of the cooling rate; if they are thermally nucleated then
$n_{0}\propto1/|r_{\mathrm{c}}|$ and $\Ton$ moves by $+5.5$~K per decade of cooling rate, by the
differentiation of Eq.~\eqref{eq:sm-logonset} carried out in Sec.~\ref{sec:unfreezing}. One
device and one afternoon separate the two.
Second, the extended area accumulated on the way down scales as $|r_{\mathrm{c}}|^{-3}$ under
continuous nucleation, so a cooldown slower than about 0.1~K/min must transform the sample on
the way down and erase the warming anomaly altogether, whichever provenance the seeds have.

On the warming branch none of this is in play. The sample begins at maximum undercooling with
the nucleus population already present, so Eq.~\eqref{eq:sm-saturated} applies, nucleation is
not rate limiting, and the only barrier left is the one that controls growth. That is why the
single-barrier description of Eq.~\eqref{eq:sm-master} is legitimate on warming and not on
cooling, and it is why the barrier extracted from the onset is a migration barrier rather than
a sum of migration and nucleation contributions.

\subsection{The isothermal hold}

Stopping the sweep is the sharpest test the mechanism admits, but the two branches test
different things and only one of them can falsify.

A hold on the cooling branch, between 150 and 170~K, interrogates the nucleation rate, because
that is the quantity the smooth cooling trace constrains. What it constrains is a lower bound.
Requiring the extended area accumulated during the sweep's passage through a given temperature
to stay below $10^{-2}$, with the Avrami exponent of three appropriate to continuous
nucleation, already places $\tau_{\mathrm{nose}}$ above 78~min at 170~K, 2.9~h at 160~K and
10~h at 150~K for a 1~K/min cooldown, and above ten to a hundred minutes at 10~K/min. A hold of
an hour that produces nothing is therefore consistent with the model as it stands; it raises the
lower bound and nothing more, and a nose time of days is not excluded by any equation here. The
cooling-branch hold is worth doing, since a hold that does transform the sample fixes
$\tau_{\mathrm{nose}}$ outright, but a null result there cannot refute the account.

The falsifying test is a hold on the warming branch, below the onset, where the nuclei are
present by hypothesis and the transformation is growth limited. The half-transformation time is
then the radius the onset demands divided by the growth velocity at the hold temperature,
$\tau_{1/2}(T)=R_{\mathrm{w}}(\Ton)/v(T)$, with both factors already fixed by the onset
analysis. Over the admissible window this is 1 to 11~s at 170~K, 5 to 49~s at 160~K, 27~s to
4.5~min at 150~K, 3 to 34~min at 140~K and 9~min to 1.9~h at 135~K, so a hold between 135 and
140~K completes in three minutes to two hours, with the central parameters giving 15 to 50~min.
That window is short enough to run and long enough to watch, and a hold there that produces
nothing is fatal to the mechanism rather than merely uninformative. The prediction has a
parameter-free anchor as well: at the onset itself the sweep has just delivered the critical
radius, so $\tau_{1/2}(\Ton)=(\kB\Ton^{2}/\Ea)/r$, which is 4 to 5~min at 1~K/min independently
of $n_{0}$, $\nu$ and $\lambda$. Anything much longer than that at $\Ton$ falsifies the model
without reference to any of its free parameters.

There is one empirical anchor for these timescales in the sister compound. Isothermal alkali
ordering in Na$_{0.8}$CoO$_{2}$ proceeds under a hold at 200~K with a time constant near 40~min
and saturates only after two to three hours, while a quench to any temperature below about 212~K
suppresses it entirely~\cite{Schulze2008}, which is the same phenomenology at the same order of
magnitude, displaced upward in temperature as the stronger alkali-alkali interaction of the
sodium system requires.

\subsection{Quasiequilibrium on the decay side}

At the upper end of the anomaly the relaxation time of the hop that carries
Eq.~\eqref{eq:sm-master},
\begin{equation}
  \tau(T)=\nu^{-1}e^{\Ea/\kB T},
  \label{eq:sm-tau}
\end{equation}
equals $1.4\times10^{-5}$~s at $T=240$~K for $\Ea=0.387$~eV and $\nu=10^{13}$~s$^{-1}$, and
$3\times10^{-4}$~s at the top of the admissible barrier range, which is shorter than a full
sweep at 1~K/min by more than seven orders of magnitude. The decay of the
anomaly is therefore a thermodynamic collapse of $\phi_{\mathrm{eq}}$ together with the shrinking
transport contrast, as in Eq.~(8) of the main text, and not a second kinetic process.

This also fixes how far the collapse can be dragged by the sweep, which is the prediction entered
in the first row of Table~II of the main text. A phase boundary crossed at rate $r$ lags its
equilibrium position by the distance the temperature moves within one relaxation time,
\begin{equation}
  \delta\Tcol \simeq \tau(\Tcol)\,r = \nu^{-1}e^{\Ea/\kB\Tcol}\,r ,
  \label{eq:sm-lag}
\end{equation}
which at $\Tcol=240$~K and the fastest rate contemplated here, $r=10$~K/min
$=0.17$~K\,s$^{-1}$, is $2\times10^{-6}$~K at the central barrier and $5\times10^{-5}$~K at
0.45~eV. The lag is exponentially sensitive to the barrier and to nothing else, so the bound is
quoted at the top of the range rather than at the middle of it: even at
0.50~eV, above the interval the onset admits, Eq.~\eqref{eq:sm-lag} gives $5\times10^{-4}$~K.
The collapse is therefore immovable below the $10^{-3}$~K level across the whole rate series and
the whole barrier range, against an onset that moves by 11~K per decade. The contrast between
Eqs.~\eqref{eq:sm-rateshift} and \eqref{eq:sm-lag} is the sharpest single statement the model
makes.

\section{Field thermodynamics of the collapse}
\label{sec:field}

\subsection{Clausius--Clapeyron relation}

Let $f_{\mathrm{P}}$ and $f_{\mathrm{O}}$ be the free energies per cobalt of the two phases in a
field $B$, with differentials $df=-S\,dT-M\,dB$, where $S$ and $M$ are the entropy and the
magnetization per cobalt. The coexistence line $\Tcol(B)$ is defined by
\begin{equation}
  \Delta f\bigl(\Tcol(B),B\bigr)=0,
  \qquad \Delta f \equiv f_{\mathrm{O}}-f_{\mathrm{P}} .
  \label{eq:sm-coex}
\end{equation}
Differentiating Eq.~\eqref{eq:sm-coex} totally with respect to $B$,
\begin{equation}
  \frac{\partial\Delta f}{\partial T}\frac{d\Tcol}{dB}
  + \frac{\partial\Delta f}{\partial B} = 0,
  \label{eq:sm-total}
\end{equation}
and using $\partial\Delta f/\partial T=-(S_{\mathrm{O}}-S_{\mathrm{P}})=\Delta S$ and
$\partial\Delta f/\partial B=-(M_{\mathrm{O}}-M_{\mathrm{P}})=-\Delta M$, with the definitions
$\Delta S\equiv S_{\mathrm{P}}-S_{\mathrm{O}}$ and $\Delta M\equiv M_{\mathrm{O}}-M_{\mathrm{P}}$
of the main text, gives Eq.~(9),
\begin{equation}
  \frac{d\Tcol}{dB}=\frac{\Delta M(B)}{\Delta S}.
  \label{eq:sm-cc}
\end{equation}
The signs are worth stating explicitly. The ordered phase has fewer accessible configurations,
so $S_{\mathrm{O}}<S_{\mathrm{P}}$ and $\Delta S>0$. The measured collapse moves upward in
field, so $d\Tcol/dB>0$, and Eq.~\eqref{eq:sm-cc} then forces $\Delta M>0$: the ordered phase is
the one with the larger magnetization. Dimensionally, magnetization per site divided by entropy
per site is (energy/field) divided by (energy/temperature), that is temperature per field, as
required.

\subsection{Numerical conversion}

Write $\Delta M=m\,\muBs$ and $\Delta S=\sigma\,\kB$ per cobalt. Then
\begin{equation}
\begin{aligned}
  \frac{d\Tcol}{dB} &= \frac{m}{\sigma}\,\frac{\muBs}{\kB},\\[2pt]
  \frac{\muBs}{\kB} &= \frac{5.788\times10^{-5}\ \mathrm{eV/T}}
                           {8.617\times10^{-5}\ \mathrm{eV/K}}
   = 0.6717\ \mathrm{K/T} .
\end{aligned}
  \label{eq:sm-convert}
\end{equation}
The measured secant slope is $\Delta\Tcol/\Delta B=(7\pm7)\,\mathrm{K}/(5\,\mathrm{T})
=1.4\pm1.4$~K/T, so
\begin{equation}
  m = \frac{1.4\pm1.4}{0.6717}\,\sigma = (2.1\pm2.1)\,\sigma ,
  \label{eq:sm-dM}
\end{equation}
which is Eq.~(10) of the main text.

Everything here is per cobalt \emph{of the phase that transforms}, and the bound below is
quoted for the reading in which that phase is the bulk. At $x\simeq1/2$ half the bulk cobalt
sites are Co$^{4+}$ in a low-spin $d^{5}$ configuration carrying $S=1/2$, so the maximum
magnetization the ordered phase can carry is $m=\msatRow$ per cobalt with every such moment
aligned, and the parent phase, being a Pauli-paramagnetic metal, contributes negligibly.
Section~\ref{sec:dftmethods} turns that ceiling from an estimate into a computed quantity, the
ground state returning \comomentRow~$\muBs$ on the Co$^{4+}$ sites and \comomentThreeRow~$\muBs$
on the Co$^{3+}$ sites. Equation~\eqref{eq:sm-dM} then bounds the latent entropy,
\begin{equation}
  \sigma \;\lesssim\; \frac{\msatRow}{2.1} \;\simeq\; \dSbound ,
  \label{eq:sm-entropybound}
\end{equation}
against the $\ln2=0.69$ per cobalt that complete lithium order and disorder at half filling would
release, so under the bulk reading the transformation converts at most about a third of the
available configurational entropy, which is consistent with the partial ordering that the
broadened rise of Sec.~\ref{sec:unfreezing} also indicates. Equation~\eqref{eq:sm-entropybound}
is not an independent statement from the headcount derived below: inverting
Eq.~\eqref{eq:sm-polar} for $\sigma$ gives
$\sigma(N)=m_{\mathrm{sat}}\tanh(N\muBs B/\kB T)/2.1$, which approaches $m_{\mathrm{sat}}/2.1$
as the block grows without limit, so $\dSbound$ is the full-polarization end of the same curve
whose other end fixes $N$. The two are an upper and a lower reading of one relation rather than
two independent results. The same calculations
show that this
ceiling is never approached in the bulk, since the dominant coupling is antiferromagnetic,
$E_{\mathrm{AFM}}-E_{\mathrm{FM}}=\eafmIntra$~meV per Co$^{4+}$ within a row, and both computed
ground states carry zero net moment, so the magnetization that tilts the collapse line has to
come from elsewhere. If it resides in a surface gas, the measured slope refers to surface cobalt
and the latent entropy per bulk cobalt is smaller by the surface-to-volume ratio, about
$10^{-2}$ for a flake 20 to 60~nm thick. Section~VII of the main text pairs calorimetry with the
thickness series to decide between the two readings.

\subsection{Induced moments}

If the phases carry no permanent moments and respond only through their susceptibilities, the
field term in the free-energy difference is $-\tfrac12\Delta\chi B^{2}$ rather than
$-\Delta M B$, and Eq.~\eqref{eq:sm-cc} integrates to
\begin{equation}
  \Tcol(B)-\Tcol(0)=\frac{\Delta\chi\,B^{2}}{2\,\Delta S}.
  \label{eq:sm-induced}
\end{equation}
The quantity extracted from a single 5~T datum is then a secant, and the terminal slope at 5~T
is twice as large, so the moment demanded at 5~T becomes $m=(4.2\pm4.2)\sigma$. The
induced-moment reading therefore requires more magnetization, not less, and strengthens rather
than weakens the headcount of the next subsection. Equation~\eqref{eq:sm-induced} also predicts
that the collapse shift grows quadratically with field, which the field series of the main text
distinguishes from the saturating shift of aligned permanent moments.

\subsection{Moment headcount}

The question the headcount answers is one of polarization rather than of energy balance. The
free-energy comparison is already exhausted by the Clausius--Clapeyron relation, which fixes the
product $\Delta M\,B$ against $\Delta S\,d\Tcol$ and says nothing about how many moments carry
the magnetization; what the measurement additionally demands is that the ordered phase actually
be magnetized by $m=2.1\,\sigma$ in units of $\muBs$ per cobalt at 5~T, and an independent spin
cannot deliver it. At $B=5$~T the single-moment Zeeman energy is
$g\muBs B=2\times5.788\times10^{-5}\times5=0.579$~meV against $\kB T$ of 12.93, 15.77 and
20.68~meV at 150, 183 and 240~K, the onset of the ordered phase, the peak of the field-doubled
excess, and the collapse, so a lone spin is outmatched by factors of 22, 27 and 36 and is
polarized only to $\tanh(\muBs B/\kB T)=0.022$, 0.018 and 0.014. The field-coupled unit must
therefore be a block of $N$ moments locked parallel by exchange, whose Zeeman energy is
$N\muBs B$ and whose polarization along the field is correspondingly larger.

Write $m_{\mathrm{sat}}$ for the magnetization per cobalt the ordered phase would carry with
every Co$^{4+}$ moment aligned, which Sec.~\ref{sec:dftmethods} computes as
$m_{\mathrm{sat}}=\msatRow$. A block of $N$ such moments is polarized to
\begin{equation}
  m = m_{\mathrm{sat}}\tanh\!\left(\frac{N\muBs B}{\kB T}\right),
  \label{eq:sm-polar}
\end{equation}
and setting this equal to the $m=2.1\,\sigma$ that Eq.~\eqref{eq:sm-dM} demands and solving for
$N$ gives Eq.~(11) of the main text,
\begin{equation}
  N = \frac{\kB T}{\muBs B}\,
  \operatorname{arctanh}\!\left(\frac{2.1\,\sigma}{m_{\mathrm{sat}}}\right),
  \label{eq:sm-headcount}
\end{equation}
which reduces to $N\simeq(2.1\sigma/m_{\mathrm{sat}})(\kB T/\muBs B)$ whenever the block is far
from saturation. The criterion carries no arbitrary constant: the latent entropy $\sigma$ enters
explicitly, so that quoting a headcount requires quoting the entropy it assumes.
Table~\ref{tab:headcount} evaluates Eq.~\eqref{eq:sm-headcount} across the entropies the
transformation admits and the three temperatures of the window.

\begin{table}[t]
  \caption{The size of the field-coupled block, Eq.~\eqref{eq:sm-headcount}, at $B=5$~T and
  $m_{\mathrm{sat}}=\msatRow$, against the latent entropy $\sigma=\Delta S/\kB$ per cobalt and
  the three temperatures at which the field effect is read. The last row is the ceiling
  $\sigma=m_{\mathrm{sat}}/2.1=\dSbound$ of Eq.~\eqref{eq:sm-entropybound}, at which the block
  is fully polarized and $N$ diverges.}
  \label{tab:headcount}
  \renewcommand{\arraystretch}{1.10}
  \begin{ruledtabular}
  \begin{tabular}{lccc}
    $\sigma$ ($\kB$/Co) & $N$ at 150~K & $N$ at 183~K & $N$ at 240~K \\
    \colrule
    0.02 & 3.9 & 4.8 & 6.3 \\
    0.05 & 9.9 & 12.1 & 15.9 \\
    0.08 & 16.3 & 19.9 & 26.1 \\
    0.10 & 21.0 & 25.6 & 33.5 \\
    0.15 & 35.1 & 42.8 & 56.2 \\
    0.20 & 60.5 & 73.8 & 96.8 \\
    \dSbound & $\infty$ & $\infty$ & $\infty$ \\
  \end{tabular}
  \end{ruledtabular}
\end{table}

At the latent entropy of about a tenth of $\kB$ per cobalt that the calorimetry row of Table~II
of the main text carries, the block holds twenty-one moments at 150~K and twenty-six at 183~K;
at half that entropy it holds ten and twelve. Across the entropies that a partial conversion
allows, 0.02 to 0.10~$\kB$ per cobalt, and across the 150 to 183~K window in which the
field-doubled excess is read, $N$ runs from about four to about twenty-six, so the field-coupled
object holds of order ten aligned moments and the range worth quoting is five to twenty-five.
Figure~3(b) of the main text draws Eq.~\eqref{eq:sm-polar} against $N$ at
150~K together with the magnetization the measured slope demands, and marks the crossings at
$N\simeq10$ for $\sigma=0.05$ and $N\simeq21$ for $\sigma=0.10$, the curve saturating at
$m_{\mathrm{sat}}$.

The energy scale required of the field-induced change in transport is a separate calculation
that happens to share a logarithm. Doubling an activated conductivity ratio requires an extra
\begin{equation}
  \delta E = \kB T\ln 2 =
  \begin{cases}
    9.0\ \mathrm{meV}, & T=150\ \mathrm{K},\\
    10.9\ \mathrm{meV}, & T=183\ \mathrm{K},
  \end{cases}
  \label{eq:sm-gapshift}
\end{equation}
of transport gap in the ordered phase, which is what Eq.~(12) of the main text asks the block of
Eq.~\eqref{eq:sm-headcount} to supply. The $\ln2$ here is the factor of two in a conductivity
ratio and nothing else; it is not a criterion for $N$, and it is unrelated to the $\ln2$ that
defines the rise midpoint in Eq.~\eqref{eq:sm-onsetdef} or to the $\ln2$ per cobalt of complete
lithium disorder in Eq.~\eqref{eq:sm-entropybound}. The three appear in three different
equations for three different reasons.

\section{Density-functional computational details}
\label{sec:dftmethods}

All calculations use \textsc{Quantum ESPRESSO} 7.3.1~\cite{Giannozzi2009,Giannozzi2017} with
the PBE exchange-correlation functional~\cite{PBE1996} and projector augmented-wave
datasets~\cite{Blochl1994} for Co, O and Li, plane-wave cutoffs of 60~Ry for the wave functions
and 480~Ry for the charge density, Marzari--Vanderbilt cold smearing~\cite{Marzari1999} of width
0.01~Ry, and a self-consistency threshold of $10^{-8}$~Ry. The Hubbard term is applied to the Co
$3d$ shell in the rotationally invariant simplified form~\cite{Dudarev1998} with atomic
projectors.

\subsection{Why $U=\Uused$~eV, and why it is not a free parameter}

The value is fixed by a physical requirement rather than chosen for agreement. A scan over
$U=0$, 2.5, 4.0 and 5.5~eV gives four qualitatively different materials. At $U=0$ every cell is
an itinerant metal with no site resolution. At $U=2.5$~eV the moment settles on the lithium-rich
cobalt rather than the lithium-poor one, so that the sign of the charge disproportionation is
inverted, and the $x=1/3$ cell is a half-metallic ferromagnet with a noninteger moment. At
$U=5.5$~eV the Co$^{4+}$ site is driven high spin, carrying \comomentHighU~$\muBs$, and the
$x=1/2$ cell is metallic again. Only near $U=\Uused$~eV does the calculation return the low-spin
$S=1/2$ Co$^{4+}$ and $S=0$ Co$^{3+}$ pair that the chemistry of Li$_x$CoO$_2$ requires. The
Hubbard parameter is therefore calibrated to a known spin state, and every number reported below
carries that calibration with it. A metal-insulator transition sits between $U=2.5$ and 4.0~eV at
$x=1/3$; above the transition the gap is nearly independent of $U$, moving from \gapXthirdUfour{} to
\gapXthirdUfive~eV between 4.0 and 5.5~eV while the moment continues to grow, which is the signature of a
gap set by the ordering pattern rather than by $U$ itself.

\subsection{Cells}

The row-ordered $x=1/2$ structure is the one-layer monoclinic cell of the Reimers--Dahn and Van
der Ven ground state~\cite{Reimers1992,VanderVen1998}, space group $P2/m$, with two Co, four O
and one Li at the experimental $a_{\mathrm{hex}}=2.808$~\AA{} and $c_{\mathrm{hex}}=14.24$~\AA.
The third lattice vector carries the O3 interlayer offset, so a single CoO$_2$ layer reproduces
the ABC stacking and the cell is a one-layer setting equivalent to the reported monoclinic cell,
with $\beta_{\mathrm{m}}=108.9^{\circ}$. The two cobalt sites are inequivalent by lithium
coordination, two neighbors against four, which permits charge order without imposing it; they are given distinct
species labels sharing one pseudopotential and one $U$, purely so that symmetrization cannot
average them. Antiferromagnetic arrangements are treated in two supercells, one doubled along
the row direction so that Co$^{4+}$ moments alternate within a row at 2.808~\AA, and one doubled
across the rows so that adjacent rows alternate at 4.864~\AA, each carrying its own ferromagnetic
and nonmagnetic references so that every comparison is made inside a single cell. Sampling is
$12\times8\times8$ for self-consistency and $16\times10\times10$ for the density of states in
the primitive cell, with equivalent densities in the supercells.

The motif comparison uses a $4\times2$ in-plane supercell, eight Co, sixteen O and four Li, that
is twenty-eight atoms and eight formula units. A $2\times2$ cell cannot do this job: on a
triangular lattice at half filling all six placements of two lithium among four sites are
nearest-neighbor pairs generating the same stripe order in one of three equivalent directions,
so row, zigzag and clumped are literally the same structure there. In the $4\times2$ cell the
row and zigzag arrangements both achieve the minimum lithium-lithium coordination of two, while
the clumped arrangement has four. Sampling is $3\times6\times4$ for self-consistency and
$4\times8\times4$ for the density of states, equivalent to $12\times12\times4$ and
$16\times16\times4$ in the primitive cell. In the row arrangement every cobalt sees
exactly three lithium, so the lithium potential is uniform over the cobalt sublattice and any
charge order there is spontaneous; the $U=0$ control confirms this, the row cell relaxing to a
uniform cobalt sublattice with moments of \momentRowUzeroA{} and \momentRowUzeroB~$\muBs$.

The $x=1/3$ calculation uses the $\sqrt3\times\sqrt3\,R30^{\circ}$ ordering with three Co, six O
and one Li per cell, sampled at $7\times7\times6$ and $12\times12\times8$.

\subsection{Caveats}

The $x=1/2$ cells are built at $c_{\mathrm{hex}}=14.24$~\AA, which rounds the 14.2235~\AA{}
interlayer spacing reported for Li$_{0.68}$CoO$_{2}$~\cite{Takahashi2007} rather than the value at
the half-lithiated composition, where the measurements give 14.42 to
14.44~\AA~\cite{Reimers1992,Motohashi2009}. The calculations
therefore sit at an interlayer distance compressed by about one percent relative to the
composition they describe, and since both the tetrahedral-site migration barrier and the
oxygen-mediated interlayer coupling are known to respond to the $c$ parameter, this is a
geometry limitation that attaches to every number computed in the row-ordered and motif cells.
No number reported here has been rescaled to compensate for it, and the comparisons that carry
the argument are made between cells sharing the same geometry, so the residual enters as a
common offset rather than as a contrast.

Hubbard projectors are atomic rather than orthogonalized atomic. Orthogonalized projectors were
tried first and were unaffordable, a single force evaluation on the fourteen-atom cell running
beyond twenty-three minutes because the derivative of the orthogonalization matrix is not
accelerated in this release, and on the ten-atom cell they also prevented self-consistency from
converging because two near-degenerate $t_{2g}$ Hubbard eigenvalues rotate into each other
between iterations. Every comparison reported here uses atomic projectors throughout, so the
choice is internally consistent, but absolute occupation traces are projector-dependent and the
charge-order claim rests on the differences between cobalt sites, which are corroborated
independently by L\"owdin $d$ occupancies and by atomic-sphere charges.

The ionic relaxation of the $x=1/2$ cell was stopped on a force plateau. The total energy became
stationary to 0.27~meV per cell over the last six steps while the reported force settled at
\forcePlateau~meV/\AA{} rather than reaching the 25~meV/\AA{} target. The residual sits entirely on the
oxygen $z$ coordinates in exact inversion pairs with the energy flat, so it is a systematic
force-expression residual rather than a real displacement, and displacements consistent with it
change the energy by less than 0.1~meV, far below the magnetic differences of tens of meV
reported here. Cobalt and lithium sit on inversion centers and carry identically zero force.

The motif comparison is made at fixed ideal geometry, identical for the three arrangements, so
that the comparison is internally consistent; the residual forces are nearly the same in all
three and are dominated by the common Co-O breathing that accompanies charge order rather than
by anything specific to the lithium configuration, so relaxation should shift the three energies
together. Ferromagnetic alignment is imposed in that campaign, and antiferromagnetic alignments
of the Co$^{4+}$ moments were not explored there.

Both the motif cells and the $x=1/3$ cell carry one CoO$_2$ layer. In the motif cells the third
lattice vector is the true O3 rhombohedral translation, so those are bulk O3 crystals rather
than slabs, but the lithium pattern is forced to repeat identically in every layer and the
specific monoclinic interlayer registry of real Li$_{0.5}$CoO$_2$ is not sampled. The $x=1/3$
cell has the exact interlayer spacing but O1 rather than O3 registry; the gap of 1.2~eV is far
larger than any plausible interlayer bandwidth and the exchange splitting is an in-plane
quantity, so neither conclusion is at risk.

Runs seeded from the high-spin start prescribed at the outset, in which
\texttt{starting\_magnetization} $=0.3$ places $5.1\,\muBs$ on every cobalt because this cobalt
dataset carries seventeen valence electrons, mostly failed to converge: seven of twenty-three
runs were abandoned after sixty to a hundred and sixteen iterations with residuals stuck between
$10^{-2}$ and $10^{-4}$~Ry, and are excluded. Where such a run stabilized enough to be read, it
sat about 1.35~eV per cobalt above the converged low-spin state, so the low-spin manifold is the
right one by a wide margin, but the high-spin branch was not exhaustively excluded at every $U$.

One methodological warning is worth recording on its own. In the seven-atom primitive cell the
ferromagnetic state comes out \primSign~meV per cobalt \emph{above} the nonmagnetic one at
$U=\Uused$~eV, which taken alone would say the ordered phase is nonmagnetic. Both supercells say
the opposite, by 26 to 36~meV per cobalt, and the nonmagnetic anchor itself agrees across all
three cells to better than 0.1~meV per cobalt. What differs is how much symmetry the magnetic
solution is permitted to break: four symmetry operations survive in the primitive cell against
two in the supercells, and the supercell solution indeed carries inequivalent Co$^{4+}$ moments
of \momentFMa{} and \momentFMb~$\muBs$, which the primitive cell cannot express. The same primitive cell
puts the magnetic state below the nonmagnetic one at every other $U$ tested. A direct control
with symmetry switched off did not converge and is reported as inconclusive, so the conclusion
rests on the agreement of the two supercells and on the cell independence of the nonmagnetic
anchor. The practical lesson holds regardless: a single primitive-cell calculation would have
returned the wrong sign for the magnetic instability in this material.

Finally, the disproportionation is about $0.2\,e$ rather than the $1\,e$ of a formal-valence
picture, the usual situation in DFT$+U$ charge-ordered oxides where the formal difference is
screened by oxygen rehybridization. The labels Co$^{3+}$ and Co$^{4+}$ describe the moment
cleanly and the charge only qualitatively. Gaps at $U=\Uused$~eV are larger than experiment
suggests for Li$_{0.5}$CoO$_2$ and should not be quoted as predictions of a gap magnitude; what survives the choice of $U$ is that the magnetic charge-ordered solutions are gapped while the
nonmagnetic ones are metallic, and that the three lithium motifs are indistinguishable, the
latter holding at $U=0$ as well.

\section{Predictions and discriminating measurements}
\label{sec:sm-predictions}

Table~\ref{tab:sm-predictions} collects the out-of-sample predictions that Sec.~VII of the main
text draws from the model, each set beside the measurement that decides it and with the numbers
and the tolerances that the derivations above supply. Every entry is out of sample in the sense
that no number in it was used in fixing the four parameters of the model, and every entry can be
performed on existing devices with existing instrumentation.

\begin{table*}[t]
  \caption{Out-of-sample predictions of the model and what each measurement decides. Rates are
  warming rates; fields are applied perpendicular to the flake unless stated.}
  \label{tab:sm-predictions}
  \renewcommand{\arraystretch}{1.15}
  \footnotesize
  \begin{ruledtabular}
  \begin{tabular}{p{0.25\textwidth}p{0.39\textwidth}p{0.28\textwidth}}
    Measurement & Prediction & What it decides \\
    \colrule
    Warming-rate series, 0.2 to 10~K/min, cooling protocol held fixed &
    Onset moves by $11\pm1$~K per decade, Eq.~(6) of the main text with the thermodynamic factor
    retained, steepening toward 16~K per decade if the cooldown is slaved to the warm-up and
    the seeds are thermally born;
    the collapse does not move, its residual kinetic lag $\tau(\Tcol)\,r$ of Eq.~\eqref{eq:sm-lag} being no
    larger than $5\times10^{-5}$~K at 10~K/min and $5\times10^{-6}$~K at 1~K/min anywhere in the
    admitted barrier range, and below $10^{-3}$~K even at 0.50~eV &
    A flat onset kills kinetic unfreezing; a collapse that follows the onset kills the
    thermodynamic assignment of $\Tcol$. The alkali order-disorder transition of the sister
    compound is sluggish enough to have been followed at more than one sweep
    rate~\cite{Ying2008} \\
    Isothermal hold at 135 to 140~K on the warming branch &
    Growth from seeds already present completes in three minutes to two hours, centrally fifteen
    to fifty minutes, after which the remaining rise is absent. A hold at 150 to 170~K on the way
    down is a different measurement, bounded only from below at one to ten hours at 1~K/min. The
    measured anchor in the sister compound is a time constant of 40~min at 200~K, with ordering
    suppressed entirely by a quench below 212~K~\cite{Schulze2008} &
    A null result on the warming branch kills the growth-from-seeds account of
    Sec.~IV\,C of the main text; a null result on the cooling branch excludes nothing \\
    Cooling-rate series, 0.05 to 10~K/min, warming rate held fixed &
    Athermal, defect-templated seeds leave the onset where it is; thermally born seeds accumulate
    as $1/|r_{\mathrm{cool}}|$ and raise the onset by 5.5~K per decade of cooling rate. A
    cooldown slower than about 0.1~K/min transforms the flake on the way down and erases the
    warming anomaly &
    Settles the provenance of the seeds, which is the one ingredient of
    Sec.~IV\,C of the main text assumed rather than derived \\
    Field series at 1, 2, 3, 5, 7, 9~T, fixed protocol &
    Saturating Brillouin growth for aligned permanent moments, $B^{2}$ growth for an induced
    moment, threshold behavior for spin-dependent percolation &
    Separates the three field couplings with no adjustable parameter \\
    In-plane against perpendicular field &
    Nearly isotropic response for spin coupling, $\cos\theta$ for orbital coupling &
    Selects spin against orbital coupling \\
    Thickness series at fixed batch &
    Excess sheet conductance proportional to thickness for a volume phase,
    thickness-independent for a surface phase &
    Decides volume against surface, and thereby whether the surface electron gas is the
    ordered phase \\
    Channel-length scaling, four-probe measurement, and a bias-history series &
    $\Delta R$ proportional to channel length for a channel origin and independent of it for a
    contact origin; amplitude independent of read voltage and of dwell time at bias if no lithium
    space charge accumulates at the electrodes &
    Confirms or kills the contact-artifact alternative of Sec.~II of the main text, including
    the Debye-thin space-charge layer that the blocking-contact argument leaves standing \\
    Batch series in $x$ at fixed protocol &
    $\Ton$ moves only through the composition dependence of $\Ea$, at $d\Ton/d\Ea=393$~K/eV, so
    the $\pm0.04$~eV environment scatter of the divacancy barrier~\cite{VanderVen2001b} allows at
    most about $\pm16$~K across $x=0.37$ to 0.48, and the same $\delta\Ea$ must appear in the
    width of the rise; $\Tcol$ may move by tens of kelvin &
    An onset shift beyond about 20~K, or one uncorrelated with the width, implicates
    inhomogeneity or a second transition rather than a single kinetic edge \\
    Calorimetry across the warming anomaly, read together with the thickness series &
    A bulk transformation gives a latent entropy of order $0.1\,\kB$ per Co, that is a latent
    heat of 2.1~meV per Co, 0.20~kJ/mol or 2.1~J/g, some twenty times a routine sensitivity of
    0.1~J/g and comparable with the 82 and 272~J/mol measured on bulk crystals at $x=0.50$ and
    $0.67$~\cite{Motohashi2009}, together with a relative excess independent of thickness; a
    surface one gives 0.02~J/g, five times below sensitivity, together with a relative excess
    inversely proportional to thickness &
    The two measurements together fix which cobalt population Eq.~(10) of the main text refers to, and
    convert it from a ratio into a measurement of $\Delta M$. No ordering exotherm on warming has
    been reported in any alkali cobaltate, so this is a new measurement rather than a
    confirmation \\
    Diffraction or phonon spectroscopy across the anomaly, set against tunneling spectroscopy &
    Three separable observables. The interlayer spacing and the $(00l)$ positions do not move.
    The in-plane metric changes by at most a tenth of a percent, with a shear below one degree,
    and the lithium-contrast superlattice reflections are weak enough to require electron
    diffraction. The site-mean Co-O bond lengths of the two cobalt sites are computed to differ
    by only \bondRow{} to \bondXthird~\AA, a magnitude that is strongly $U$-sensitive and is not
    itself a prediction of the method, while a hard gap opens in the local spectrum &
    Separates electronic charge order from a displacive or Jahn-Teller transition, which would
    contrast the two sites by tens of picometers, and from a stacking glide, which would collapse
    the interlayer spacing by six percent \\
    Electron or synchrotron diffraction on one flake at 100~K and at 300~K &
    Weak lithium-contrast superlattice intensity present in the warmed-through low-temperature
    state and absent in the as-cooled state, with the interlayer spacing and the $(00l)$
    positions unchanged to better than half a percent throughout &
    Confirms that what orders is the lithium sublattice, and excludes the O3 to H1-3 stacking
    glide of Sec.~II\,A of the main text, which collapses the interlayer spacing by about six
    percent \\
  \end{tabular}
  \end{ruledtabular}
\end{table*}

\section{Symbols}
\label{sec:symbols}

Table~\ref{tab:symbols} collects every symbol that carries physical content in the main text and
in this Supplemental Material, together with its dimension, so that a symbol appearing in a
caption and the same symbol appearing inside a derivation can be checked against one another.

\begin{table*}[t]
  \caption{The symbols that carry physical content in the main text and in this Supplemental
  Material, with their dimensions. The left panel collects the transport and rate quantities,
  the right panel the nucleation, thermodynamic and electronic-structure quantities. Purely
  local auxiliaries introduced and discharged inside a single derivation, such as the
  integration variable $y$ of Eq.~\eqref{eq:sm-beta} and the substitutions $\varepsilon$,
  $a_{0}$ and $C$ of Eq.~\eqref{eq:sm-sensitivity}, are defined where they appear and are not
  repeated here. No symbol in the list below carries a second meaning anywhere in either
  document.}
  \label{tab:symbols}
  \renewcommand{\arraystretch}{1.10}
  \footnotesize
  \begin{ruledtabular}
  \begin{tabular}{lll@{\qquad}lll}
    Symbol & Meaning & Dimension & Symbol & Meaning & Dimension \\
    \colrule
    $\phi$ & transformed volume fraction & dimensionless & $I_{\mathrm{n}}$, $I_{0}$ & nucleation rate per area and its prefactor & m$^{-2}$s$^{-1}$ \\
    $\phi_{\mathrm{eq}}$ & equilibrium transformed fraction & dimensionless & $v$, $v_{0}$ & domain growth velocity and its prefactor & m\,s$^{-1}$ \\
    $x$ & lithium content of Li$_{x}$CoO$_{2}$ & dimensionless & $n_{0}$ & frozen nucleus areal density & m$^{-2}$ \\
    $q$ & thermodynamic factor $1-e^{\Delta f/\kB T}$ & dimensionless & $a$ & lithium-lithium spacing in the plane & m \\
    $s$ & resistivity ratio $\rho_{\mathrm{O}}/\rho_{\mathrm{P}}$ & dimensionless & $\lambda$ & domain-wall line tension & eV\,m$^{-1}$ \\
    $c$ & mixture geometry factor & dimensionless & $W^{*}$, $R^{*}$ & nucleation barrier and critical radius & eV, m \\
    $R_{\mathrm{meas}}$ & measured two-probe resistance & $\Omega$ & $A_{\mathrm{ex}}$ & extended area fraction & dimensionless \\
    $R_{\mathrm{P}}$ & channel resistance of the parent phase & $\Omega$ & $A_{0}$, $n$ & Avrami prefactor and exponent & dimensionless \\
    $R_{\mathrm{c}}$, $R_{\mathrm{ch}}$ & contact and channel resistance & $\Omega$ & $\Delta f$ & free-energy difference per cobalt & eV \\
    $\Delta R$ & warming-minus-cooling excess & $\Omega$ & $\Delta f_{A}$ & free-energy difference per unit area & eV\,m$^{-2}$ \\
    $L$, $A$ & channel length and cross section & m, m$^{2}$ & $\Delta S$ & latent entropy per cobalt & eV\,K$^{-1}$ \\
    $k(T)$ & rearrangement rate & s$^{-1}$ & $\Delta S_{A}$ & latent entropy per unit area & eV\,K$^{-1}$m$^{-2}$ \\
    $\nu$ & attempt frequency & s$^{-1}$ & $\sigma$ & $\Delta S$ in units of $\kB$ & dimensionless \\
    $\Ea$ & effective migration barrier & eV & $\Delta M$ & magnetization difference per cobalt & eV\,T$^{-1}$ \\
    $\delta\Ea$ & spread of barriers across the sample & eV & $m$ & $\Delta M$ in units of $\muBs$ & dimensionless \\
    $r$ & sweep rate & K\,s$^{-1}$ & $\Delta\chi$ & susceptibility difference per cobalt & eV\,T$^{-2}$ \\
    $T_{0}$ & temperature at which the sweep starts & K & $B$ & applied magnetic field & T \\
    $\Theta$ & accumulated attempts on the sweep & dimensionless & $g$ & spectroscopic splitting factor & dimensionless \\
    $\delta T_{10\times}$ & temperature per decade of $A_{\mathrm{ex}}$ & K & $N$ & moments in the field-coupled block & dimensionless \\
    $\Delta T_{10\text{--}90}$ & 10 to 90 percent rise width & K & $\delta E$ & field-induced change of transport gap & eV \\
    $R_{\mathrm{w}}$ & growth radius accumulated on warming & m & $m_{\mathrm{sat}}$ & $\Delta M$ with every Co$^{4+}$ moment aligned & dimensionless \\
    $\Ton$ & rise midpoint of the anomaly & K & $U$ & Hubbard parameter on Co $3d$ & eV \\
    $\Tcol$ & first-order collapse temperature & K & $D(E_{F})$ & density of states at the Fermi level & eV$^{-1}$ \\
    $\delta\Tcol$ & kinetic lag of the collapse & K & $n_{d}$ & cobalt $3d$ occupation & dimensionless \\
    $\tau$ & relaxation time $\nu^{-1}e^{\Ea/\kB T}$ & s & $\beta_{\mathrm{m}}$ & monoclinic cell angle & degrees \\
    $\tau_{\mathrm{nose}}$ & transformation time at the nose & s & $\mathcal{H}$ & thermal history of the sample & dimensionless \\
    $\tau_{1/2}$ & isothermal half-transformation time & s & $X$ & areal density of nucleation templates & m$^{-2}$ \\
    $\psi$ & heterogeneous nucleation potency & dimensionless & & & \\
  \end{tabular}
  \end{ruledtabular}
\end{table*}

\clearpage

\bibliographystyle{apsrev4-2}
\bibliography{refs}

\end{document}
